\section{Thiết lập thực nghiệm}\label{sec:setup}

\paragraph{Dữ liệu.}
Chúng tôi dùng LiHua-World, bộ dữ liệu do chính nhóm tác giả MiniRAG xây dựng~\cite{fan2025minirag}: 442 tài liệu hội thoại tiếng Anh của một người dùng giả định, mỗi tài liệu mang mốc thời gian.
Bộ câu hỏi có 637 dòng (635 câu phân biệt) chia ba loại: Single (một bước), Multi (nhiều bước) và Null (đáp án đúng là ``không đủ thông tin'').
Chúng tôi tách 200 câu làm tập dev, phân tầng theo loại với seed cố định.
Mọi chẩn đoán, sàng lọc và lựa chọn cấu hình đều làm trên tập dev.
\textbf{Kết quả chính được đo trên \nHeld{} câu ngoài dev} (\nHeldSingle{} Single, \nHeldMulti{} Multi, \nHeldNull{} Null). Tập này không được dùng để chọn biến thể nào.

\paragraph{Đối chứng.}
Đối chứng chính là MiniRAG tái lập, chạy cùng điều kiện với mọi biến thể (\cref{sec:fixes}).
Vectơ thuần là đối chứng đơn giản nhất mà bài này đòi hỏi.
Số của bài MiniRAG gốc chỉ là \emph{tham chiếu}: bài gốc dùng giám khảo khác và một đồ thị khác. Chúng tôi không tính hiệu số với nó.

\paragraph{Độ đo.}
Mỗi câu trả lời được giám khảo gán một trong ba nhãn: \emph{accurate} (truyền đạt đúng đáp án vàng), \emph{error} (khẳng định điều mâu thuẫn với đáp án vàng mà không rào đón) và \emph{neither} (từ chối, không biết hoặc lạc đề).
Chúng tôi báo tỉ lệ của cả ba nhãn (acc, err, neither) và acc/(acc+err).
Tỉ số cuối tách được hiệu ứng ``dám trả lời nhiều hơn'' khỏi hiệu ứng ``trả lời đúng hơn''.
Với câu Null, lời từ chối có nghĩa là trả lời đúng.

\paragraph{Giám khảo.}
Giám khảo là \texttt{gemini-flash-lite-latest} với prompt ở \cref{app:prompts}.
Mỗi lượt sinh dùng lượt chấm thứ nhất. Một câu được tính là đúng khi phiếu \emph{accurate} chiếm đa số tuyệt đối có trọng số; hai dòng câu hỏi trùng văn bản được gộp, mỗi dòng nặng một nửa.
Trên MiniRAG tái lập, độ lệch chuẩn accuracy giữa ba lượt chấm của cùng một lượt sinh là \sdJudgeBase{} điểm.
Độ tin cậy của giám khảo so với người được bàn ở \cref{sec:limitations}.

\paragraph{Chi tiết cài đặt.}
Mô hình sinh là \texttt{Qwen/Qwen2.5-3B-Instruct}~\cite{qwen2025qwen25}, phục vụ bằng vLLM 0.11.0 với cửa sổ 32.768 token. Mô hình này dùng cho cả lập chỉ mục, phân tích câu hỏi lẫn sinh câu trả lời.
Nhiệt độ lấy mẫu không được đặt trong yêu cầu, nên dùng giá trị mặc định của máy chủ \todo{xác nhận nhiệt độ thực tế}.
Vectơ nhúng là \texttt{all-MiniLM-L6-v2}~\cite{reimers2019sbert,wang2020minilm} (384 chiều), giống bài gốc.
Đoạn dài tối đa 1.200 token; đồ thị do Qwen2.5-3B dựng có \graphNodes{} nút và \graphEdges{} cạnh.
Tham số truy vấn: \texttt{top\_k} = 60, $\budget = 4000$ token cho phần đoạn.
Mọi cấu hình dùng cùng một chỉ mục đã đóng băng.

\paragraph{Lượt chạy.}
Nhiễu giữa các lượt sinh lớn hơn nhiễu giám khảo: trên \nHeld{} câu, ba lượt sinh của RRF(đồ thị, vectơ) có độ lệch chuẩn \sdGenGV{} điểm.
Vì vậy mỗi biến thể được sinh ba lượt. Các lượt chạy cho đăng ký trước dùng seed cố định $\{101, 202, 303\}$; ba lượt của RRF(đồ thị, vectơ) và lượt đầu của vectơ thuần có từ trước, không đặt seed, và được khoá làm mốc trước khi chạy BM25. Không chọn seed tốt nhất.
Riêng MiniRAG tái lập chỉ có \emph{một} lượt sinh; hệ quả của điều này được nêu ở từng phép so liên quan.

\paragraph{Kiểm định thống kê.}
Kiểm định chính là hoán vị đổi dấu ghép cặp trên hiệu accuracy trung bình theo câu, với 100.000 hoán vị và seed 20260914.
Khoảng tin cậy 95\% là bootstrap ghép cặp theo câu (10.000 lần lấy mẫu lại).
Ba giả thuyết sau được \textbf{đăng ký trước}, trước khi viết mã BM25 và trước mọi đầu ra; chúng được hiệu chỉnh Holm~\cite{holm1979} trên ba phép so:
\begin{itemize}
  \item \textbf{H1}: RRF(đồ thị, vectơ, BM25) so với RRF(đồ thị, vectơ) --- thêm BM25 khi có đồ thị.
  \item \textbf{H2}: RRF(vectơ, BM25) so với vectơ thuần --- thêm BM25 khi không có đồ thị.
  \item \textbf{H3}: RRF(vectơ, BM25) so với RRF(đồ thị, vectơ, BM25) --- bỏ đồ thị khi đã có vectơ và BM25.
\end{itemize}
Một giả thuyết chỉ được gọi là \emph{đạt} khi thoả cùng lúc các điều kiện sau: hiệu trung bình $\ge T = \Tthr$ điểm (hai lần độ lệch chuẩn giữa các lượt sinh); cả ba cặp lượt cùng chiều dương; $p_{\text{Holm}} < 0{,}05$; err không tăng quá 1 điểm; không nhóm câu nào giảm có ý nghĩa và nhóm Null không mất trung bình quá 3 câu; token ngữ cảnh nằm trong $\pm 5\%$ của 3.870.
Bốn phép so với MiniRAG và giữa vectơ thuần với RRF(đồ thị, vectơ) \textbf{không} nằm trong đăng ký trước. Chúng tôi báo chúng như phân tích khám phá, hiệu chỉnh Holm riêng trên bốn phép so.
Từng cặp lượt được báo kèm bằng số câu lên/xuống (sai$\to$đúng / đúng$\to$sai).
Trong toàn dự án chúng tôi đã thử 23 cấu hình hoặc hướng (\cref{app:tried}).
